\begin{figure}[H]
\centering
\resizebox{\textwidth}{!}{%
\begin{tikzpicture}[
  font=\small\sffamily,
  layer/.style={
    draw=VIUDark!80,
    fill=VIULight,
    rounded corners=3pt,
    align=center,
    text width=3.15cm,
    minimum height=1.05cm,
    inner sep=5pt
  },
  market/.style={
    draw=VIUGray!70,
    fill=VIUWhite,
    rounded corners=3pt,
    align=center,
    text width=3.35cm,
    minimum height=1.05cm,
    inner sep=5pt
  },
  gap/.style={
    draw=VIUOrange,
    fill=VIUOrange!10,
    rounded corners=4pt,
    align=center,
    text width=4.25cm,
    minimum height=1.28cm,
    inner sep=7pt,
    font=\small\sffamily\bfseries
  },
  claim/.style={
    draw=VIUDark,
    fill=VIUDark!5,
    rounded corners=4pt,
    align=center,
    text width=4.2cm,
    minimum height=1.28cm,
    inner sep=7pt
  },
  arrow/.style={-{Latex[length=2.2mm]},thick,color=VIUDark!85},
  oarrow/.style={-{Latex[length=2.2mm]},thick,color=VIUOrange!95!black},
  faint/.style={-{Latex[length=2mm]},thick,color=VIUGray!70}
]
\node[market] (m1) at (0,3.0) {Mercado AMR\\fleet manager, WMS, tráfico};
\node[market] (m2) at (3.9,3.0) {Material handling\\picking, pallets, tugging};
\node[market] (m3) at (7.8,3.0) {Carga compatible\\un robot o plataforma dedicada};

\node[layer] (a1) at (0,1.05) {MRTA y mercados\\subastas, CBBA, coaliciones};
\node[layer] (a2) at (3.9,1.05) {Grafos y consenso\\vecinos, DAC, conectividad};
\node[layer] (a3) at (7.8,1.05) {Control físico\\formación, MPC, CBF, wrench};

\node[layer] (l1) at (0,-0.9) {MARL/GNN\\políticas, message passing};
\node[layer] (l2) at (3.9,-0.9) {Geometría\\rigidez, contactos, slots};
\node[layer] (l3) at (7.8,-0.9) {Robustez\\batería, fallos, comunicación};

\node[gap] (gap) at (11.85,1.05) {Brecha focal\\asignación distribuida + quorum entero + factibilidad física};
\node[claim] (smith) at (11.85,-0.9) {Smith-QR\\payoff de déficit, cierre QR, capacidad efectiva y guardias de ejecución};

\draw[faint] (m1) -- (m2);
\draw[faint] (m2) -- (m3);
\draw[arrow] (a1) -- (a2);
\draw[arrow] (a2) -- (a3);
\draw[arrow] (l1) -- (l2);
\draw[arrow] (l2) -- (l3);

\draw[oarrow] (m3.east) -- (gap.west);
\draw[oarrow] (a3.east) -- (gap.west);
\draw[oarrow] (l3.east) -- (gap.west);
\draw[oarrow] (gap.south) -- (smith.north);

\node[draw=VIUGray!60,fill=VIUWhite,rounded corners=3pt,align=center,
  text width=8.35cm,inner sep=5pt,font=\scriptsize\sffamily] (note) at (3.9,-2.45)
  {Lectura: el mercado ha madurado en orquestación de flotas y transporte unitario; la academia cubre asignación, control, aprendizaje y física. El hueco defendible está en cerrar esas capas en una ley distribuida auditable para una misma carga compartida.};
\draw[faint] (l2.south) -- (note.north);
\end{tikzpicture}%
}
\caption[Taxonomía del estado del arte para transporte cooperativo multi-robot]{Taxonomía del estado del arte para transporte cooperativo multi-robot de cargas. La contribución se ubica en el puente entre asignación distribuida y factibilidad física de una carga compartida.}
\label{fig:sota-transporte-cooperativo}
\end{figure}
